\section{Setting and notation}
\label{sec:preliminaries}

We work in a deterministic, finite-dimensional, single-token setting: there is no
randomness and no asymptotics, only inner-product geometry together with a first-order
optimality condition. All objects below are fixed throughout.

\paragraph{Carrier and vocabulary.}
Let $E$ be a real inner product space with inner product $\inner{\cdot}{\cdot}$ and induced
norm $\norm{\cdot}$; we assume $E$ is complete (e.g. $E = \mathbb{R}^{d}$ with the Euclidean
inner product, or any real Hilbert space). Let $\Vocab$ be a finite vocabulary with at least
two tokens, and let
\begin{align*}
W : \Vocab \to E,
\end{align*}
be the unembedding map, so that $W_{a} \in E$ is the row associated with token $a \in \Vocab$.

\paragraph{LayerNorm sphere.}
Fix a radius $r > 0$. The LayerNorm sphere of radius $r$ is
\begin{align*}
\Sphere \;=\; \{\, x \in E : \norm{x} = r \,\}.
\end{align*}
A point $x$ lies on the sphere precisely when $\norm{x} = r$.

\begin{definition}[Decoder: greedy argmax is the true decoder]
\label{def:decoder}
Let $\loss : E \to \mathbb{R}$ denote the implicit loss and let $x \in E$ be a hidden state.
Because the setting is single-token ($n = 1$), the greedy argmax over rows \emph{is} the true
decoder. We say that token $a \in \Vocab$ is \emph{generated at} $x$, written
$\Generated(a, x)$, when $x$ lies in the open decoding cone of $a$, i.e. the row $W_{a}$ has
strictly larger inner product with $x$ than every competing row:
\begin{align*}
\Generated(a, x) \;\;:\Longleftrightarrow\;\; \forall\, b \ne a,\quad \inner{W_{a}}{x} > \inner{W_{b}}{x}.
\end{align*}
\end{definition}

\begin{definition}[Constrained stationarity on the sphere]
\label{def:stationarity}
We model ``the hidden state has settled at a constrained stationary point of the loss on the
sphere'' by the first-order (Lagrange) condition: a point $\xstar \in \Sphere$ is
\emph{constrained-stationary} with multiplier $\mu \in \mathbb{R}$ when the ambient gradient
of the loss is normal to the sphere at $\xstar$,
\begin{align*}
\grad \loss(\xstar) \;=\; \mu\,\xstar.
\end{align*}
Equivalently, the Riemannian gradient of $\loss$ along the sphere vanishes at $\xstar$. We
take the multiplier $\mu$ and this stationarity equation as hypotheses; no convexity, global
minimality, or convergence rate is assumed.
\end{definition}

\begin{remark}[Differentiability is a setting remark, not a hypothesis]
\label{rem:differentiability}
The loss $\loss$ is intended to be differentiable, and $\grad\loss(\xstar)$ denotes its
gradient at $\xstar$. This is a property of the ambient setting, not a working hypothesis of
the results below: the proofs never differentiate $\loss$; they use only the \emph{value} of
the vector $\grad\loss(\xstar)$ through the stationarity equation of
\Cref{def:stationarity}. Accordingly, differentiability of $\loss$ is deliberately omitted
from the statements, which keeps them as general as the argument actually warrants.
\end{remark}
